\section*{INTRODUCCIÓN}
\addcontentsline{toc}{section}{INTRODUCCIÓN}
\parrafoI{
    La integración de la Inteligencia Artificial Generativa y los Modelos de Lenguaje de Gran Escala ha provocado un cambio de paradigma sin precedentes en la ingeniería de software. Estas herramientas han demostrado una notable capacidad para acelerar la producción automatizada de código, trasladando progresivamente la lógica de desarrollo tradicional hacia la formulación de instrucciones semánticas o prompts. Sin embargo, la adopción acelerada de esta tecnología ha revelado una carencia crítica: la ausencia de marcos metodológicos formales que garanticen la coherencia estructural y la sostenibilidad del software producido.  

En la práctica actual, el uso empírico de la IAG suele priorizar la funcionalidad inmediata sobre la calidad interna. Esta tendencia genera sistemas con una elevada complejidad ciclomática y deriva en la acumulación de lo que la literatura reciente denomina ``PromptDeb'' o deuda de prompt. Este fenómeno compromete directamente los atributos de mantenibilidad exigidos por estándares internacionales como la norma ISO/IEC 25010, propiciando la creación de sistemas opacos o ``cajas negras'' que resultan insostenibles técnica y económicamente a mediano y largo plazo. Esta problemática es especialmente visible a nivel nacional y regional, donde la transformación digital y la adopción tecnológica en entornos académicos y empresariales, como se observa en la región Cusco, requieren urgentemente de estándares de gobernanza técnica rigurosos.  

Frente a este escenario, la presente investigación tiene como propósito fundamental evaluar el efecto de la implementación de un Framework de Ingeniería de Prompts, sustentado en los lineamientos de calidad de la norma ISO/IEC 25010, sobre la mantenibilidad y la deuda técnica del código generado por inteligencia artificial. A través de un diseño cuasi-experimental, el estudio propone la aplicación sistemática de patrones arquitectónicos de prompts y la integración de pipelines automatizados de análisis estático, buscando transformar el proceso estocástico de generación de código en una disciplina de ingeniería predictiva, medible y auditable.
}